\section{Justificación de la investigación}

\subsection*{Justificación teórica}
El tema abordado es de plena actualidad: el paradigma RAG fue formalizado apenas en 2020 y la integración de mecanismos de autorreflexión en agentes LLM introduciendo modelos críticos es aún más reciente (2023-2025), con un volumen de publicaciones en crecimiento sostenido según se evidenció en el Estudio de Mapeo Sistemático conducido como fase previa de este trabajo. La investigación aportada por la sistematización del conocimiento sobre la intersección de RAG y autorreflexión, un espacio que según se evidencia en el SMS conducido por {\cite{Alcides2026}}; evidenciando un campo insuficientemente explorado por la literatura existente, concentrada mayormente en modelos de lenguaje de gran escala o de dominio general. 

Si bien \textcite{Asai2024_Selfrag} con su herramienta Self-RAG; resuelven la construcción del componente crítico mediante un modelo entrenado por separado y supervisado con anotaciones obtenidas de un modelo propietario, que inserta offline los tokens de reflexión sobre el corpus de entrenamiento del generador para Dominio general, deja una brecha abierta en sus resultados promoviendo la adopción o trabajos futuros a dominios específicos. \textcite{jeong_improving_2024}, publicado en \textit{Bioinformatics} (ISMB 2024), se enfoca puntualmente sobre Dominios específicos en Biomedicina con Self-BioRAG tomando como base Self-RAG. Los autores afirman que al aplicar los métodos existentes a distintos problemas de dominio específico se hace evidente una generalización deficiente, que conduce a recuperar documentos incorrectos o a emitir juicios inexactos. La doble consecuencia: no solo falla la recuperación (documentos incorrectos), sino también el juicio (evaluaciones inexactas). La calidad de ese modelo crítico se valida por su concordancia (superior al 90 \% en la mayoría de las categorías de tokens) con las predicciones de referencia del modelo propietario; se trata, por tanto, de una validación dentro de la distribución de datos sobre la que el crítico fue entrenado (Dominio específico).

El trabajo contribuye conceptualmente a establecer las relaciones y compromisos entre confiabilidad del razonamiento, capacidad paramétrica y restricciones de recursos computacionales. La salida que la literatura ofrece a este problema: reentrenar el componente crítico con instrucciones del dominio específico, anotadas mediante llamadas a un modelo propietario. De esta doble incompatibilidad emerge la pregunta que el Modelo Critic-EmpRAG1.0 busca responder: cómo dotar de contexto de dominio específico al componente crítico sin reentrenarlo sobre datos del dominio propietario, anclando su juicio en la evidencia recuperada en tiempo de inferencia o utilizando un Modelo Híbrido para la Autorreflexión.

\subsection*{Justificación práctica}

La propuesta del Modelo crítico,  Critic-EmpRAG1.0, persigue un propósito transformador de corto plazo sobre la realidad de Dualbiz SRL: la mejora de la productividad y el acceso al Sistema de soporte de conocimiento empresarial interno de las áreas de TI y desarrollo. Esto se traduce en tiempos de ejecución de proyectos más cortos, menor tiempo dedicado a la gestión de soporte y una disminución del costo de oportunidad asociado a los retrasos actuales. De no realizarse esta investigación, la organización continuaría dependiendo de mecanismos manuales de transferencia de conocimiento, con los consiguientes sobrecostos y retrasos descritos en los antecedentes. De realizarse, se espera una reducción medible de dichos tiempos y una mayor autonomía del personal en la resolución de consultas operativas.

Sin acción de por medio, la situación problemática identificada tiende a agravarse con el tiempo: a medida que Dualbiz SRL. y organizaciones en circunstancias similares continúan creciendo, el volumen de conocimiento operativo interno y la cantidad de personal que requiere acceder a él aumentan, mientras los mecanismos manuales de transferencia de conocimiento no escalan al mismo ritmo. La propuesta es además viable y factible: se apoya en modelos de lenguaje pequeños de código abierto, frameworks de RAG y técnicas de autorreflexión ya validados individualmente en la literatura (\cite{Shinn2023_Reflexion}; \cite{Asai2024_Selfrag}), y cuenta con un caso de estudio real y accesible para el investigador (Dualbiz SRL) para su validación empírica.



\subsection*{Justificación social y económica}

Desde un enfoque social, de acuerdo con  \textcite{Achury2016induccion}, dotar al personal de Dualbiz SRL. tanto nuevo como existente de un mecanismo confiable de acceso al conocimiento interno reduce la dependencia de supervisión intensiva por parte de personal experimentado. Por otra parte, agiliza la incorporación de nuevos colaboradores y mejora la calidad de las decisiones cotidianas con base en dicho conocimiento. Desde el punto de vista económico, una inducción más eficiente reduce costos de capacitación y permite a la empresa asumir nuevos proyectos con mayor rapidez, mientras que sus clientes se benefician de tiempos de respuesta de soporte más cortos. Adicionalmente, al tratarse de una arquitectura pensada para operar bajo restricciones de \textit{hardware} y presupuesto, sus resultados resultan potencialmente extensibles, sujeto a validación posterior, a otras pequeñas y medianas empresas de \textit{software} con necesidades similares, ampliando su relevancia social más allá del caso de estudio particular.

\subsection*{Justificación metodológica}

Desde lo metodológico, la investigación enriquece el procedimiento que se propone para diseñar y validar arquitecturas de autorreflexión sobre pipelines RAG, específicamente la elaboración del modelo crítico, Critic-EmpRAG1.0, y el desempeño en dominio específico (empresarial); utilizando modelos de lenguaje pequeños, que puedan ser más eficientes que los modelos generalizados. Un procedimiento que parte de un mapeo sistemático de la literatura (SMS) para fundamentar las decisiones de diseño que puede ser reutilizado por otros investigadores u organizaciones que enfrenten el mismo problema en contextos empresariales con restricciones en infraestructura de computo y limitación de acceso con privacidad en la información, similares a Dualbiz SRL. y que puedan ampliar a otros dominios específicos la arquitectura planteada para los modelos Críticos híbridos.

